\chapter*{Vorwort}
\addcontentsline{toc}{chapter}{Vorwort}

\begin{merksatz}
\boxtitel{Die Lücke, die dieser Band schließt}
Band 4 behandelte den Zufall in einem Augenblick: eine Ziehung, eine Verteilung, ein Erwartungswert. Die Anwendungsbände brauchten aber immer wieder etwas anderes --- Zufall, der \kursiv{weitergeht}.

\begin{center}
{\small
\renewcommand{\arraystretch}{1.5}
\begin{tabular}{@{}P{2.2cm}P{6.4cm}@{}}
\toprule
\textbf{Band 16} & Irrfahrten auf Graphen, ohne ihre Theorie.\\
\textbf{Band 18} & Epidemien als Sprungprozesse, näherungsweise durch Differentialgleichungen ersetzt.\\
\textbf{Band 22} & Gradientenabstieg mit Zufallsschritten.\\
\textbf{Band 24} & Markow-Entscheidungsprozesse --- Prozesse im Namen, ohne Prozesstheorie.\\
\bottomrule
\end{tabular}}
\end{center}

\begin{center}
\begin{minipage}{0.9\textwidth}
\textbf{Fünfmal wurde derselbe Begriff gebraucht und jedes Mal nur so weit entwickelt, wie es gerade nötig war. Dieser Band holt ihn nach.}
\end{minipage}
\end{center}
\end{merksatz}

\begin{merksatz}
\boxtitel{Drei Ergebnisse, die man nicht erwartet}
\begin{center}
{\small
\renewcommand{\arraystretch}{1.55}
\begin{tabular}{@{}P{1.0cm}P{7.4cm}@{}}
\toprule
$1$ & Bei einem fairen Spiel führt \kursiv{nicht} meist einer knapp --- am wahrscheinlichsten ist, dass einer fast die ganze Zeit vorn liegt. \kursiv{(Kapitel 1)}\\
$2$ & Eine Irrfahrt in der Ebene kehrt mit Sicherheit zum Ausgangspunkt zurück, im Raum nur mit Wahrscheinlichkeit $0{,}34$. \kursiv{(Kapitel 1)}\\
$3$ & Die Brownsche Bewegung ist überall stetig und nirgends differenzierbar --- und genau deshalb braucht man ein eigenes Integral, dessen Kettenregel einen Zusatzterm trägt. \kursiv{(Kapitel 10--12)}\\
\bottomrule
\end{tabular}}
\end{center}

\begin{center}
\begin{minipage}{0.88\textwidth}
\kursiv{Das erste Ergebnis widerspricht der Anschauung so gründlich, dass man es mehrfach nachrechnen möchte. \textbf{Es besagt, dass in einer Serie fairer Spiele die längste Glückssträhne die Regel ist und nicht die Ausnahme} --- und es erklärt, warum Menschen in zufälligen Folgen unentwegt Muster sehen.}
\end{minipage}
\end{center}
\end{merksatz}

\begin{merksatz}
\boxtitel{Die sechs Teile}
\begin{center}
{\small
\renewcommand{\arraystretch}{1.5}
\begin{tabular}{@{}P{1.4cm}P{3.6cm}P{3.0cm}@{}}
\toprule
\textbf{Teil} & \textbf{Thema} & \textbf{Kapitel}\\
\midrule
I & Ketten & $1$--$3$\\
II & Sprünge in stetiger Zeit & $4$--$6$\\
III & Faire Spiele & $7$--$9$\\
IV & Die Brownsche Bewegung & $10$--$12$\\
V & Anwendungen & $13$--$14$\\
VI & Stand & $15$\\
\bottomrule
\end{tabular}}
\end{center}

\begin{center}
\begin{minipage}{0.88\textwidth}
\kursiv{\textbf{Die Teile I bis III kommen mit Band 4 aus} --- Summen, Erwartungswerte, bedingte Wahrscheinlichkeiten. Erst Teil~IV verlangt die Maßtheorie aus Band 13, und auch dort nur ihre Grundbegriffe. Wer die Brownsche Bewegung zunächst überspringt, kann Teil~V bis auf die Optionsrechnung lesen.}
\end{minipage}
\end{center}
\end{merksatz}

\begin{warnung}
\boxtitelr{Drei Warnungen}
\begin{center}
{\small
\renewcommand{\arraystretch}{1.55}
\begin{tabular}{@{}P{2.2cm}P{6.4cm}@{}}
\toprule
\textbf{Erste} & Die Anschauung trügt bei Prozessen stärker als irgendwo sonst in der Mathematik. Fast jedes Ergebnis dieses Bandes verletzt eine plausible Erwartung.\\
\textbf{Zweite} & Ein Modell, das einen Verlauf gut beschreibt, muss ihn nicht erzeugt haben. Dieselben Pfade entstehen aus ganz verschiedenen Mechanismen.\\
\textbf{Dritte} & Die Formeln der Finanzmathematik in Kapitel~13 setzen Annahmen voraus, die an den Rändern regelmäßig brechen --- und die Ränder sind das, worauf es dort ankommt.\\
\bottomrule
\end{tabular}}
\end{center}

\begin{center}
\begin{minipage}{0.88\textwidth}
\kursiv{Die zweite Warnung verbindet diesen Band mit Band 23. \textbf{Aus einem beobachteten Verlauf auf den erzeugenden Mechanismus zu schließen, ist ein kausales Problem} --- und die Prozesstheorie liefert dafür Modelle, keine Entscheidung zwischen ihnen.}
\end{minipage}
\end{center}
\end{warnung}

\begin{merksatz}
\boxtitel{Was dieser Band voraussetzt}
\begin{center}
{\small
\renewcommand{\arraystretch}{1.48}
\begin{tabular}{@{}P{3.4cm}P{1.4cm}P{2.8cm}@{}}
\toprule
\textbf{Werkzeug} & \textbf{Band} & \textbf{gebraucht in}\\
\midrule
Wahrscheinlichkeit & $4$ & durchgehend\\
Matrizen, Eigenwerte & $2$ & Kapitel 2--3\\
Grenzwerte, Reihen & $3$ & Kapitel 1--6\\
Maß und Integral & $13$ & ab Kapitel 7\\
Differential\-gleichungen & $14$ & Kapitel 5, 12\\
\bottomrule
\end{tabular}}
\end{center}
\end{merksatz}
